\begin{table*}[t]
\centering
\caption{Selected corrosion evidence for AM HEAs in chloride media. All entries are laboratory saline tests unless noted; long-term natural seawater data for AM HEAs remain unavailable.}
\label{tab:corrosion}
\footnotesize
\begin{tabularx}{\textwidth}{@{}llllX@{}}
\toprule
Alloy & AM route & Medium & Key finding & Source \\
\midrule
CoCrFeMnNi & LPBF & 3.5\% NaCl & Higher pitting resistance and polarization resistance than cast; $\sim$58\% lower corrosion current density vs.\ cast & \cite{han2025,cagirici2024} \\
CoCrFeMnNi & LPBF (as-built, HIPed) & 3.5\% NaCl vs.\ H$_2$SO$_4$ & $\sim$2.2$\times$ better than 316L in chloride; $\sim$15.5$\times$ worse in acid; low H permeation vs.\ 316L & \cite{lpbf_cantor_2024} \\
AlCoCrFeNi$_{2.1}$ & SLM & 3.5\% NaCl & As-built homogeneous duplex outperforms annealed and cast; annealing restores segregation and degrades passivity & \cite{slm_ehea_2025} \\
AlCoCrFeNi$_{2.1}$Mo$_x$ & (cast; AM-adjacent design) & 3.5\% NaCl & Mo ($x{=}0.3$) raises strength and pitting potential via Mo/Cr composite oxide film & \cite{mo_ehea_2026} \\
AlMo$_{0.25}$FeCoCrNi$_{2.1}$ & LDED & Saline & Annealing heat treatment jointly tunes strength and corrosion response & \cite{lded_almofeco2024} \\
AlCoCrFeNi (Ti, Cu doped) & LMD & Saline & Doping improves corrosion resistance over undoped; intermetallics must be suppressed & \cite{cagirici2024} \\
Nb--Mo--Ta--W & LPBF & Saline & Better overall than 316L; segregation sites produce pits up to $\sim$1~mm deep & \cite{cagirici2024} \\
CoCrW & LPBF & Acidic NaCl, 28~d immersion & Oxide film matures then destabilizes; attack at melt-pool boundaries & \cite{lpbf_coocrw2023} \\
\bottomrule
\end{tabularx}
\end{table*}
